\chapter[Pleasure, Virtue, or Tranquillity?]{Is Happiness Pleasure, Virtue, or Freedom from Disturbance?}
\section{A Discarded Wooden Cup}
Begin with an object: a wooden drinking cup.

No gold-rimmed museum antique, just hollowed and polished wood used daily to scoop water. Over twenty-three centuries ago in Corinth, it belonged to an old man nicknamed the Dog. Diogenes slept and ate in the street without caring what others thought. People called him a Cynic, dog-like, an insult that later named a philosophical school.

Diogenes Laertius's \textit{Lives of Eminent Philosophers} preserves a familiar story. Seeing a child drink from cupped hands, Diogenes paused, removed his wooden cup from his bag, and threw it away. A child had surpassed him in learning to live: even a cup was unnecessary baggage. Later he saw a child break bread by hand without a spoon and discarded his spoon too.

Pause over the strangeness. We normally discard things because we obtain something better: an old cup for a new one. He discarded his because he discovered no cup was needed. His possessions dwindled to a cloak, staff, and begging bag; supposedly even the bag nearly went.

Put that discarded cup on your desk today. What lies there? Phone, headphones, water bottle, snacks, revision guides, a photograph of the new trainers you want. Our generation owns more than any ancient king, yet the feeling of still wanting something scarcely diminishes. Advertising always points towards one more thing. Diogenes faces the opposite way and asks through his discarded cup: what if happiness lies towards one fewer?

This is no exhortation to hardship but an entrance to a real question. Is happiness pleasurable feeling, upright living, or undisturbed peace? Greeks seriously argued for over two centuries and produced answers still fighting within our lives. Let us listen.

\section{Athenian Sunshine and an Emperor's Shadow}
Come into a scene.

A sunny morning in the fourth century BCE, in Corinth's square. Mediterranean light dazzles white; paving stones grow hot. Olive oil, salted fish, and dried figs scent the air. Vendors call, while strolling citizens in robes debate yesterday's politics beneath colonnades.

Beside the square stands an enormous ceramic storage jar, once used for grain or wine, now inhabited. Diogenes makes it his house. Wrapped in an old cloak, he lies in sunlight, eyes shut, like a warmed stone.

A disturbance rolls down the street. Crowds part, vendors fall silent, philosophers stop arguing. Alexander approaches, no ordinary man but Alexander the Great, conqueror of much of the known world, in his twenties and Greece's most powerful person. Curious about the jar-dwelling eccentric, he has come himself.

His large shadow falls across Diogenes. According to the \textit{Lives}, he announces himself as Alexander, the great king. The man beside the jar answers, ``I am Diogenes, the Dog.'' Alexander offers any favour he might desire.

Humanity has remembered the next reply for over twenty-three centuries. Broadly:

``Then move aside, and stop blocking my sunlight.''

Feel its weight. This visitor can grant gold, palaces, or office with one sentence, or order death. The man on the ground inventories what he already has---sunshine, warmth, free time---and declares every possible gift inferior to the sunlight behind the king. All worldly power amounts here to one nuisance: obstructed light.

Legend says Alexander was not angry, telling attendants that were he not Alexander, he would choose to be Diogenes. Its truth is untestable, but generations repeat it because it voices a hidden unease: is the world-owning conqueror or the possessionless sunbather nearer happiness?

\section{What Are We Asking? A Misframed Question}
Neither applaud Diogenes nor dismiss him as a performance artist yet. State the conflict.

We often hear ``Are you happy?'' in street interviews, school essays, and companies' annual reviews. It hides a trap: happiness is assumed to be one measurable thing, like blood pressure, higher or lower, yielding an immediate answer.

Greeks thought otherwise. Their word, \textit{eudaimonia}, usually translated as happiness, is nothing like feeling cheerful. \textit{Eu} means good, \textit{daimon} originally a guarding spirit or deity. Together, roughly a life watched over by a good spirit: living well, having a good life. Not feeling good now but living well overall. This distinction grounds the chapter.

Momentary pleasure resembles sugar: immediately, certainly, briefly sweet. Greek happiness resembles a tree. One tasted leaf cannot judge it; inspect its soil, storms, and eventual growth. A tree's goodness takes time to see; a good human life almost requires a lifetime's retrospect.

Now comes a real conflict, not mere wording.

One side speaks from bodily intuition: happiness is pleasure, surely? Good exam results and holidays feel genuinely good. If that is not happiness, what is? Why should philosophers dismiss real feelings?

The square's thinkers stand opposite. Aristotle says pleasure is one component, not the whole. Epicurus says it is indeed the goal, but everyone misunderstands genuine pleasure. Stoics say neither pleasure nor pain should rule you; inner peace should. Diogenes's discarded cup declares the supposed necessities themselves the disease.

Four schools use eudaimonia for almost four different lives. No translation accident, but an unfinished human battle. Are ``a good university and job bring happiness'' and ``enjoyment matters most; ease the pressure'' not fighting in your home now?

Before continuing, answer honestly: forced to choose, would you prefer a pleasant but ordinary life, or a difficult but worthwhile one? Remember and compare after reading.

\section{Four Answers, Four Lives}
Invite each disputant to speak. None asks how to enjoy today alone, but what makes an entire life good.

\textbf{Aristotle: Happiness means playing your hand well, throughout a lifetime.}

Aristotle taught Alexander. Yes, the emperor and the eccentric connect through the same philosopher. In the \textit{Nicomachean Ethics}, he proposes an enduring idea: everything has a distinctive function and excellence. Eyes see, instruments make music. A thing's goodness lies in performing its function well. A good knife cuts sharply; a good musician plays excellently.

What about humans? Their distinctive capacity is rational action. Living well means exercising it excellently throughout life, an excellence he calls virtue, \textit{arete}. This can mislead if understood narrowly as moral niceness. It means excellence: courage facing fear, moderation facing desire, wisdom in judgement. Virtue is no inborn certificate but a muscle trained like sport or music. Repeated courageous acts make a courageous person.

Two points follow. Happiness is activity, not state: not a possession like a good mood or mark, but the life being lived. And it requires a lifetime's account, not one exam or holiday. Aristotle even asks whether posthumous disgrace or descendants' misfortune changes the goodness of a previously happy life. Happiness is scarcely settled before the coffin closes. He is no pretender to lofty detachment either: health, friendship, some wealth, and reasonable fortune are conditions. Without them, virtue struggles to flower.

\textbf{Epicurus: Pleasure is the answer, but you have it backwards.}

Epicurus may be history's most misunderstood philosopher. English epicurean suggests gourmet indulgence, and Chinese pleasure-seeking often implies dissipation. Reality is nearly opposite. Pause at this easy misjudgement: the doctrine named after him is precisely not his doctrine.

Pleasure is indeed his goal, but its greatest form is absence of pain: bodily health and freedom from anxiety. He bought a garden outside Athens and lived with friends almost austerely on bread and water, an occasional cheese making a feast. His reasoning is cool: indulgence brings bills---hangovers, vanity's emptiness, competitive troubles. Total the account and it loses. Satisfy natural necessary desires, eating when hungry and drinking when thirsty; discount empty wants for fame, display, luxury; then evade the two great fears, divine punishment and death.

His secret weapon is friendship, astonishingly high on his happiness list. Knowing someone will come when trouble strikes supplies security, among pleasure's steadiest sources. Unlike the Cynic who cuts away needs, Epicurus cultivates garden, table, and friends.

\textbf{Stoics: Distinguish what you control.}

Their name comes from an Athenian painted colonnade, \textit{stoa}, where they first taught. This school lasted longest and travelled farthest, Athens to Rome, followed by enslaved people and emperors. Its central idea: some things depend on you, others do not. Judgements, choices, responses do; exam difficulty, others' liking, rain, relatives' illness do not. Much suffering comes from confusing columns, trying desperately to control the second while neglecting what the first permits.

Their happiness centres on peace and freedom from disturbance, stability that external events cannot overturn. They speak of following nature or cosmic order: the world follows its workings, and wisdom treats facts as questions to answer, not opponents in a tug-of-war. Calling Stoicism expressionless emotional repression is another error. They deny not love or grief, but that happiness must await permission from uncontrollable things.

\textbf{Cynics: Tear up society's list.}

Return to Diogenes's jar. Cynics are most radical: almost every civilised desire---wealth, status, reputation, respectability---is an unnecessary chain. People spend lives chasing chains and becoming possessions' slaves. Live naturally: take only what is needed, refuse the rest, including others' opinions. Public coarseness, jar-dwelling, not rising before a king all demonstrate freedom from convention. Alone among the four, they overturn even the agreement that good living requires some money.

Together they form a spectrum: Aristotle, live well; Epicurus, calculate clearly; Stoics, remain steady; Diogenes, cut away. They share a premise unfamiliar today: happiness concerns living an entire life before it concerns rating this moment's mood.

Nor was the dispute exclusively Greek. Around the same time, Confucius praised his favourite pupil Yan Hui in ``Yong Ye'' in the \textit{Analects}:

\begin{quote}
A basket of food, a gourd of water, in a poor lane: others could not bear its hardship, yet Hui never altered his joy.
\end{quote}
Plain rice, water, a shabby lane: unbearable to others, still joyful for Yan Hui. Confucius was similar. ``Transmitting'' records:

\begin{quote}
Eating coarse food, drinking water, resting my head on a bent arm: joy is there too. Wealth and honour gained unjustly are drifting clouds to me.
\end{quote}
Coarse grain, plain water, an arm as pillow contain pleasure; unjust riches and rank resemble clouds. Without exchanging letters, Greece and China reached strikingly similar structures: pleasure need not come from external conditions. India's Buddha approached from another direction, locating suffering in wanting and holding tight, and release in understanding and loosening attachment. Three unconnected traditions asked who holds happiness's sovereignty, the world or oneself.

Before applauding detachment, honestly strengthen the other side: everyone seriously pursuing pleasure, money, and comfort, perhaps you and your parents. Their case is powerful. First, pleasure and pain are real. Asking a hungry child to discuss spiritual peace is cruel hypocrisy: food before philosophy. Second, those saying money cannot buy happiness often have enough; for poor families, extra money supplies genuine happiness, amply supported by experience. Third, pursuit of material improvement advances society. If everyone merely basked, who would build irrigation, cure disease, establish schools? Fourth, lowering desires easily becomes the powerful telling the poor to stay quiet. ``Be content'' means something different from a boss than from a friend.

Do not dismiss this. Most philosophical answers address those whose basic needs are met but who remain anxious about something else. The hungry and the fed do not face the same happiness question.

\section[Thinking Bridge: Three Happiness Questions]{Thinking Bridge: Divide Happiness into Three Questions}
After the dispute, take away a tool. Next time someone asks whether you are happy, first ask which kind they mean.

Divide happiness into three questions. Many disputes turn out to involve two people answering three different questions.

\textbf{First: How do you feel now? (Emotion)}

Joy or sadness, relaxation or tension: minutes and hours. Did you enjoy today? This layer is real, direct, fleeting. A bad mood may ruin an afternoon but usually not a year. Psychologists can measure it through ratings throughout the day.

\textbf{Second: Are you satisfied with life overall? (Evaluation)}

Step back over recent life---study, friends, family, health---and rate it across months or years. It separates from the first layer. Irritation at a problem can coexist with a fulfilling term; daily fun can end in yearly emptiness. Poor emotion, high evaluation, or the reverse.

\textbf{Third: Is your life worth living? (A whole life)}

This is the Greek concern: neither feeling nor satisfaction scores, but who you are becoming and whether your deeds are worthwhile, measured over a lifetime. Often it demands present unpleasantness: practising music, distance running, caring for illness, losing out to honour a promise. The first account loses; the third may gain. Aristotle almost entirely occupies this layer.

Apply it at home. Mum says less gaming and better marks bring future happiness. You answer that life is short and enjoyment matters most. Apparently rival happiness views, actually Mum addresses lifetime worth and you immediate pleasure, neither first checking the question. The tool does not resolve disagreement, but locates it: not whether to seek happiness, but which layer takes priority.

Use the same blade on advertising. ``Buy it; you deserve a better self.'' Advertising switches between the second and third layers, packaging shopping anticipation, the first, as an upgraded life, the third. Recognise the switch and half consumer rhetoric becomes clear.

The tool ends here. It identifies the question, not which matters most. That decision belongs to values, not technique.

\section{Reading an Ancient Letter and Handbook}
Read brief primary material to hear the thinkers themselves.

First, Epicurus's letter to his friend Menoeceus, copied into the \textit{Lives of Eminent Philosophers}, among the most direct evidence of his thought. A frequently quoted passage says, broadly:

\begin{quote}
When we call pleasure the goal, we mean neither the pleasures of dissipation nor sensual indulgence, as those who misunderstand us suppose. We mean freedom from bodily pain and disturbance of the soul.
\end{quote}
Notice its form: a private letter, not a paper, like a long chat message today. Its content almost corrects a rumour. Already called dissipated, he insists that his pleasure means not hurting and not panicking. A philosopher's crucial work becomes preventing misunderstanding, which survives two millennia. That itself shows how much people need pleasure-seeking to mean indulgence before comfortably despising it.

Second, the later Stoic teacher Epictetus. His life comments upon his philosophy: born enslaved in Rome's world, speaking of inner freedom. That is the point, not irony. A pupil compiled his teaching into the \textit{Handbook}, whose first sentence says, broadly:

\begin{quote}
Some things depend on us, others do not. Judgement, impulse, desire, aversion---our own actions---depend on us. Body, property, reputation, position---everything not our own action---do not.
\end{quote}
Feel its weight from an enslaved person. His body legally belonged to someone else; reportedly his leg was disabled too. External conditions compressed the first two happiness layers severely. He built his philosophy on the third: others can imprison the body, not how he regards it. That final inch remains his. Disagree, but do not call him frivolous. His life tested the theory.

Together these passages reveal Hellenistic philosophy's medicinal posture. Epicurus explicitly calls it treatment for fear and greed. Epictetus makes a portable tool to consult when events strike: is this mine to control? If not, why should it overturn me? This classical page belongs beside today's books on managing emotion and ending mental exhaustion, only nearly two thousand years earlier.

We quote them not because ancient people were always right, but because late-night advice to lighten up has traceable roots. Knowing them helps distinguish grounded wisdom from inspirational broth brewed for clicks.

\section[Was Diogenes Happy? Three Explanations]{Was Diogenes Actually Happy? Three Explanations}
Compare explanations of one fact: Diogenes abandoned possessions, home, respectability, yet claimed happiness, believed by many contemporaries, while empire-owning Alexander supposedly envied him. Why? Each account has reasons and limits.

\textbf{Explanation One: Subtraction as treatment---deleting suffering's background programs.}

Epicurean and Cynic reasoning suggests an algorithm: suffering equals what one wants minus what one owns. Most enlarge the numerator; Diogenes cuts the denominator nearly to zero. Own nothing, lose nothing; expect nobody's help, suffer nobody's disappointment. ``Stop blocking the sun'' is not pretended superiority but an almost desireless person's honest financial report: the account currently balances. Internally coherent and partly familiar---think of relief after clearing a to-do list. He clears life's entire list. Limits are firm: can desires vanish? He must beg, feels cold, falls ill. Minimal survival does not mean zero suffering. And can one person's solution scale to a city? If everyone sleeps in jars, who farms or heals?

\textbf{Explanation Two: Performance---happiness as art before an Athenian audience.}

Change position. Every gesture is public: square, jar, crowd. Cup-discarding awaits witnesses, royal defiance circulation. Here he does not leave society but depends intensely upon it. His happiness project requires audiences, discussion, generations of teachers retelling him, including this book now. His product is not happiness but \textbf{criticism}: his body becomes a placard against civilisation, demonstrating that its coveted goods are unnecessary. This explains theatricality and why cynic now suggests knowing, sneering disillusionment in Chinese and English. Its limit: performance cannot be falsified; any sustained austerity can be called acting. It also cannot explain followers genuinely living this way for centuries. Pure theatre cannot last that long.

\textbf{Explanation Three: Defensive response---guarding the final inch in an uncontrollable world.}

Place him historically. Greek cities declined through incessant war, then Alexander's conquests swept away the old world. Yesterday a respected citizen, today your city vanished. External happiness became hazardous, so all four schools moved its foundations from world to mind. The jar becomes a lucid defensive structure, not madness: when everything can be seized, keep only the unseizable. Stoicism's later appeal to both enslaved people and emperors follows similarly; those unable to control fate find it persuasive. This restores philosophy to history and explains its timing. Its limit: reducing thought to psychological massage for an age's illness, a losers' consolation, when it also works in peace and among successful people.

Each holds part of the truth. Perhaps their overlap comes closest: amid an uncontrolled age, someone sincerely practises subtraction through performance. Methodologically, whether he really was happy is evidence, whether his approach is replicable is strategy, and whom it criticises is politics. Understanding requires separating questions entangled in one person.

A Chinese parallel: around the same time, ``Autumn Floods'' in the \textit{Zhuangzi} tells of Chu's king inviting Zhuangzi to office. Fishing in the Pu River, he replies without turning. A sacred Chu tortoise, dead three thousand years, rests honoured in a shrine. Would it prefer honoured death or living with its tail in mud? Living, say the envoys. Then leave, says Zhuangzi; I too intend to trail my tail in mud. Mediterranean sunshine and Pu River mud answer one question across a continent.

\section[Further Challenge: Measuring Happiness]{Further Challenge: Can Happiness Be Measured?}
Now the weightiest theoretical problem. Greeks asked what happiness is; our age adds \textbf{whether it can be measured}, disputed across ethics, psychology, and economics.

Begin with measurers. Psychology takes happiness into laboratories through surprisingly simple methods: satisfaction questionnaires or devices asking randomly about current feelings. Decades yield reasonably robust findings. Hedonic adaptation: excitement over good things fades, new-phone joy in weeks, moving home or improved marks in months, before returning to a baseline and wanting again. Hence the hedonic treadmill: running without changing scenery. Money and happiness also form a steep-then-flattening curve. Below basic provision, money matters greatly; beyond a threshold, added happiness diminishes. This supports rather than contradicts our material realists: money matters for the poor; another purchase does not guarantee happiness for the comfortable. More surprisingly, people predict happiness badly. They expect exam failure to hurt for a year but recover in days, lottery wins to delight for life but adapt quickly.

These findings arm ancient arguments. Adaptation supports Epicurus: chasing stimulation loses value as pleasure dilutes. Diminishing returns beyond a money threshold hand Aristotle the microphone: above it, relationships, engagement, meaning, and autonomy matter.

But ethicists identify a deep difficulty: \textbf{questionnaires may measure something other than the Greek question}. The three-layer tool makes this plain. Scales mostly address present emotion and life satisfaction. Aristotelian eudaimonia concerns a virtuous life's unfolding; how do you score yourself on that? Caring for a gravely ill relative may bring daily negative feelings and barely passing satisfaction, yet be the life's least-regretted period. A scale records unhappiness; the person objects. Who is right?

Psychologist and economics Nobel laureate Daniel Kahneman describes two selves: an experiencing self living each moment and bearing emotions, and a remembering self narrating afterwards. Their accounts diverge. A mostly delightful holiday with a terrible final day becomes a bad memory; an arduous mountain ascent causes suffering throughout but becomes a brightest-life entry. In our terms, experience inhabits Layer One, memory Layer Three, while ``Are you happy?'' never specifies whom it asks.

The dispute's structure is now visible. Measurers say improvement requires measurement; policy cannot rest on poets' adjectives. Opponents say compressing happiness into a score quietly answers the ancient philosophical question before measuring. Happiness is \textbf{defined} as what scores capture; dignity, meaning, and virtue slipping off the ruler are declared nonexistent. Both are partly right. A ruler can compare which life feels better, not which is more worthwhile. The latter requires your judgement.

\section[Dopamine, Marks, and Emotional Stability]{Today: Dopamine, Report Cards, and Emotional Stability}
The ancient fight continues in phones and classrooms, each school with a modern double.

Epicurus's difficulty is industrialised pleasure. Videos, games, snacks are carefully tuned pleasure extractors, more precise than ancient feasts. Adaptation runs at full power: a stimulus every fifteen seconds, a faster treadmill, until normal-speed lessons, books, and walks become intolerable. He would note the irony: these products are anti-Epicurean, supplying empty unnecessary wants and sending bills for attention, sleep, and emptiness afterwards. His old accounting still concludes a loss overall.

Stoicism has revived fully. Emotional stability is prized; distinguishing controllable from uncontrollable appears unchanged in popular psychology and self-help. Cognitive behavioural therapy, a leading treatment for anxiety, shares Epictetus's skeleton: interpretations, not events alone, trouble us. Exam anxiety supplies practice. Paper difficulty, others' marks, cut-off scores occupy Column Two. Tonight's revision and use of two examination hours occupy One. No inspirational platitude, but direct application of the columns, and genuinely effective.

Aristotle lives in long-term thinking and deliberate practice. Virtue as trained habit means you do not become disciplined before acting so; repeated disciplined action makes you disciplined. Countless habit books repeat this, seldom crediting the fourth century BCE.

The Cynics' descendants are most intriguing. Everyday cynicism means disbelief, meaninglessness, and sneering withdrawal, almost opposite Diogenes. His refusal is passionate and committed; modern refusal often coldly abandons commitment. A travelling concept reaches its origin's opposite over two thousand years. Beware another misjudgement: a classmate apparently detached from everything may be tired, disappointed, or protecting themselves, not Diogenes. Ancient Cynicism attacks actively in sunshine; it does not hide under a duvet declaring the world pointless.

A final contemporary situation: elders and teachers speak Layer Three---be useful, do worthwhile things---while peers comfort one another in Layer One---enjoy yourself, stop exhausting yourself mentally. This is not generations at war but three questions arguing through two generations' mouths. Identifying the question is a practical contribution to the dinner table.

\section[Your Turn: The Endless-Pleasure Machine]{Your Turn: The Machine of Unending Pleasure}
End with a question without a standard answer.

Twentieth-century philosopher Robert Nozick proposed a famous experiment. Imagine scientists build a machine that connects to your brain and supplies uninterrupted, utterly convincing pleasurable experiences: success, love, adventure, glory, everything desired, forever. One cost: your body floats in a tank, nothing happens outside the machine, and you never know.

First question: is a life connected to it happy?

If you recoil despite flawless feelings, Aristotle lives in you. You care not only for experience but whether things happened and you truly lived. If you would connect, judging pleasure sufficient and reality's distinction mere attachment, you side coherently with strict hedonism.

Second, nearer life: many choices are smaller machines. A night scrolling beats a night practising music on Layer One. Telling everyone what they want to hear wins more comfort than truthfulness risking isolation. Repeated choices establish an exchange rate between present good feeling and intangible worthwhileness.

What is your rate? More importantly, \textbf{who sets it}? Parents, algorithms, or you? Two centuries of Greek argument concerned this sovereignty. Aristotle assigns it to the self shaped daily; Epicurus to the clear accountant; Stoics to the inch you control. Diogenes says first tear up the list others thrust upon you, then discuss it.

They reached no consensus, nor have two millennia since. Perhaps this is not a problem awaiting a final standard answer but a carving tool accompanying each whole life. Your answers carve who you become.

Give your answer and your reasons.
